%%%%%%%%%%%%%%%%%%%%%%%%%%%%%%%%%%%%%%%%%%%%%%%%%%%%%%%%%%%%%%%%%%%%%%
% Lớp 10 | Đại số tổ hợp | Nhị thức Newton
% Dạng 0D8?3-3: Tìm hệ  số, số hạng trong khai triển bằng tam giác Pascal
% Số câu: 8  (sinh tự động bởi tach_cau.py; từ nay thầy sửa trực tiếp file này)
%%%%%%%%%%%%%%%%%%%%%%%%%%%%%%%%%%%%%%%%%%%%%%%%%%%%%%%%%%%%%%%%%%%%%%

%%%%%%%%%%%%%%%% TRẮC NGHIỆM 4 PHƯƠNG ÁN %%%%%%%%%%%%%%%%

% nguồn: MaHoaNH_TracNghiem_Lop10
\begin{ex}%[0D8H3-3]%[ma:0c69a137]
	Hệ  số của số hạng thứ chính giữa trong khai triển $\left(x^2-\dfrac{2}{x}\right)^{10}$ là
	 \choice
		 {$-\mathrm{C}_{10}^52^5x^5$}
		{$\mathrm{C}_{10}^52^5x^5$} 
		 {$\mathrm{C}_{10}^52^5$}
		 {\True $-\mathrm{C}_{10}^52^5$}
	\loigiai{ 
		Ta  có $\left(x^2-\dfrac{2}{x}\right)^{10}=\sum\limits\limits_{k=0}^{10} \mathrm{C}_{10}^k(x^2)^{10-k} \cdot\left(\dfrac{-2}{x}\right)^k$.\\
		Số hạng chính giữa ứng với  $k=5$. Khi đó hệ số của nó là $\mathrm{C}_{10}^5\cdot(-2)^5=-\mathrm{C}_{10}^52^5$.
	 }
\end{ex}

% nguồn: MaHoaNH_TracNghiem_Lop10
\begin{ex}%[0D8H3-3]%[ma:c2ee67dc]
	Cho nhị thức $(2x+3y)^n$. Tìm $n$ biết hệ số của số hạng thứ $3$  chia cho số hạng thứ $2$ trong khai triển theo số mũ giảm dần của $x$ bằng $\dfrac{9}{2}$.
	\choice 
		{$n=5$} 
		{$n=6$} 
		 {\True $n=7$}
		 {$n=8$}
	 \loigiai{
		Ta có $(2x+3y)^n=\sum\limits_{k=0}^n  \mathrm{C}_n^k(2x)^{n-k} \cdot(3y)^k$.\\
		Hệ số của số hạng thứ 2 ứng với $k=1$ là $P_2=\mathrm{C}_n^12^{n-1} \cdot  3^1$.\\
		Hệ số của số hạng thứ 3 ứng với $k=2$  là $P_3=\mathrm{C}_n^22^{n-2} \cdot 3^2$.\\
		Suy ra $\dfrac{\mathrm{C}_n^22^{n-2} \cdot 3^2}{\mathrm{C}_n^12^{n-1} \cdot 3^1}=\dfrac{9}{2} \Leftrightarrow \dfrac{n(n-1) \cdot 2^{n-2} \cdot 3^2}{2n \cdot 2^{n-1}  \cdot 3}=\dfrac{9}{2} \Leftrightarrow n-1=6\Leftrightarrow n=7$.
	} 
\end{ex}

% nguồn: MaHoaNH_TracNghiem_Lop10
\begin{ex}%[0D8H3-3]%[ma:b2cf1c55]
	Tìm số hạng không chứa $x$ trong khai triển của nhị thức $\left(x^3-\dfrac{1}{x^2}\right)^5$. 
	 \choice
		 {\True $-10$}
		{$-5$} 
		 {$10$}
		 {$5$}
	 \loigiai{
		Ta có 
		 $$
		\begin{aligned}
			\left(x^3-\dfrac{1}{x^2}\right)^5 &=(x^3)^5-5(x^3)^4\left(\dfrac{1}{x^2}\right)+10(x^3)^3\left(\dfrac{1}{x^2}\right)^2-10(x^3)^2\left(\dfrac{1}{x^2}\right)^3+5(x^3)\left(\dfrac{1}{x^2}\right)^4-\left(\dfrac{1}{x^2}\right)^5  \\
			&=x^{15}-5 x^{10}+10  x^5-10+5\cdot \dfrac{1}{x^5}-\dfrac{1}{x^{10}}.
		\end{aligned}
		 $$
		Số hạng không chứa $x$ trong khai triển là $-10$. 
	 }
\end{ex}

% nguồn: MaHoaNH_TracNghiem_Lop10
\begin{ex}%[0D8H3-3]%[ma:74181aa9]
	Số hạng chứa $\sqrt{x}$ trong khai triển $\left(\sqrt{x}+\dfrac{2}{x}\right)^4$, $x > 0$ là số hạng thứ  mấy?
	 \choice
		 {$5$}
		 {$3$}
		 {\True $2$}
		{$4$} 
	 \loigiai{
		Ta có 
		\allowdisplaybreaks 
		\begin{align*}
			\left(\sqrt{x}+\dfrac{2}{x}\right)^4&=\left(\sqrt{x}\right)^4+4\left(\sqrt{x}\right)^3\left(\dfrac{2}{x}\right)+6\left(\sqrt{x}\right)^2\left(\dfrac{2}{x}\right)^2+4\left(\sqrt{x}\right)\left(\dfrac{2}{x}\right)^3+\left(\dfrac{2}{x}\right)^4 
			 \\
			&=x^2+8\sqrt{x}+24\cdot \dfrac{1}{x}+32\cdot  \dfrac{\sqrt{x}}{x^3}+16\cdot \dfrac{1}{x^4}.
		\end{align*}
		 Số hạng chứa $\sqrt{x}$ trong khai triển trên ứng với số hạng thứ $2$.
	} 
\end{ex}

% nguồn: MaHoaNH_TracNghiem_Lop10
\begin{ex}%[0D8H3-3]%[ma:8e4a8223]
	Biết $\left(1+\sqrt[3]2\right)^4=a_0+a_1  \sqrt[3]2+a_2 \sqrt[3]4$. Tính $a_1a_2$.
	\choice 
		{$a_1 a_2=24$} 
		{$a_1  a_2=8$}
		{$a_1 a_2=54$} 
		 {$a_1 a_2=36$}
	\loigiai{ 
		Ta  có
		 \allowdisplaybreaks
		\begin{align*}
			\left(1+\sqrt[3]{2}\right)^4&=1^4+4\cdot  1^3\left(\sqrt[3]{2}\right)^1+6\cdot 1^2\left(\sqrt[3]{2}\right)^2+4\cdot 1^1\left(\sqrt[3]{2}\right)^3+\left(\sqrt[3]{2}\right)^4
			 \\
			&=1+4 \sqrt[3]{2}+6  \sqrt[3]{4}+8+2 \sqrt[3]{2} \\
			&=9+6 \sqrt[3]{2}+6 \sqrt[3]{4}. 
		\end{align*}
		 Suy ra $a_1\cdot a_2=6\cdot 6=36$.
	} 
\end{ex}

% nguồn: MaHoaNH_TracNghiem_Lop10
\begin{ex}%[0D8H3-3]%[ma:06aa54e3]
	 Số hạng chứa $x^4$ trong khai triển đa thức $\left(2x^2+3x+1\right)(2x+1)^4$ là
	 \choice
		{\True  $160x^{4}$}
		 {$160$}
		 {$80x^{4}$}
		{$80$} 
	\loigiai{ 
		Ta có  $\left(2x^2+3x+1\right)(2x+1)^4=(x+1)(2x+1)(2x+1)^4=(x+1)(2x+1)^5$.
		 \allowdisplaybreaks
		\begin{align*}
			(x+1)(2x+1)^5&=(x+1)  \sum\limits\limits_{i=0}^5\mathrm{C}_5^i(2x)^{5-i} 1^i
			\\ 
			&=(x+1) \sum\limits_{i=0}^5\mathrm{C}_5^i 2^{5-i} x^{5-i} 
			\\ 
			&=\sum\limits\limits_{i=0}^5\mathrm{C}_5^i 2^{5-i} x^{6-i}+\sum\limits_{i=0}^5\mathrm{C}_5^i 2^{5-i}  x^{5-i}.
		\end{align*}
		Số hạng chứa $x^4$ trong khai  triển là $\mathrm{C}_5^22^3x^4+\mathrm{C}_5^12^4x^4=160x^4$.
	 }
\end{ex}


%%%%%%%%%%%%%%%% ĐÚNG/SAI %%%%%%%%%%%%%%%%

% nguồn: MaHoaNH_DungSai_Lop10
\begin{ex}%[0D8H3-3]%[ma:0727caaf]
	$\mathrm{A}^k_n$; $\mathrm{C}^k_n$; $\mathrm{P}_n$ lần lượt là số chỉnh hợp chập $k$ của $n$ phần tử, số tổ hợp chập $k$ của $n$ phần tử và số hoán vị của $k$ phần tử. Trong các khẳng định sau, khẳng định nào \textbf{đúng}? 
	\choiceTF 
		 {\True $\mathrm{P}_n=n!$}
		 {$\mathrm{C}^{k-1}_n +\mathrm{C}^k_n =\mathrm{C}^{k+1}_{n+1}$}
		 {\True $\mathrm{C}^k_n=\mathrm{C}^{n-k}_n$}
		{$\mathrm{A}^k_n=\dfrac{\mathrm{C}^k_n}{k!}$} 
	 \loigiai{Ta có
		\begin{eqnarray*}
			&&\mathrm{P}_n =n!.\\ 
			 &&\mathrm{A}^k_n=\dfrac{n!}{(n-k)!}.\\
			&&\mathrm{C}^k_n=\dfrac{n!}{(n-k)!k!}=\dfrac{\mathrm{A}^k_n}{\mathrm{P}_k.}\\ 
			&& \textbf{[Công  thức Pascal]} \mathrm{C}^{k-1}_{n-1} +\mathrm{C}^k_{n-1} =\mathrm{C}^k_{n}.
	\end{eqnarray*}}
\end{ex}

% nguồn: MaHoaNH_DungSai_Lop10
\begin{ex}%[0D8H3-3]%[ma:3b94859d]
	Sử dụng nhị thức Newton, khai triển $\left(3x+2\right)^4$ và tìm hệ thức đúng trong các hệ thức đúng. 
	 \choiceTF
		{\True Hệ số  của số hạng chứa $x^2$ là $216$}
		 {Hệ số của số hạng chứa $x$ là $24$}
		{\True Hệ số của số hạng chứa $x^3$ là  $108$}
		{Hệ số của số hạng chứa $x^4$ là  $16$}
	 \loigiai{
		Ta có 
		\begin{eqnarray*}
			\left(3x+2\right)^4 &=&(3x)^4+4\cdot(3x)^3\cdot2+6\cdot(3x)^2\cdot2^2+4\cdot(3x)\cdot2^3+2^4\\ 
			&=&81x^4+108x^3+216x^2+96x+16. 
		\end{eqnarray*}
		Vậy hệ số của $x^3$ là $108$ và hệ  số của số hạng chứa $x^2$ là $216$.
	 }
\end{ex}
